\begin{figure}[ht]
\centering
\begin{tikzpicture}[x=1cm, y=1cm, font=\small]

% Three stacked horizontal intervals sharing a common centre, showing
% the width ordering CI < PI < TI for the same sample. Centre at x=5.
\def\cx{5.0}

% Common centre line
\draw[dashed, enggray] (\cx, -0.4) -- (\cx, 3.6);
\node[below, font=\scriptsize, enggray] at (\cx, -0.4) {$\bar{x}$};

% Confidence interval (narrowest), y = 3
\draw[very thick, engblue] (\cx-0.7, 3) -- (\cx+0.7, 3);
\draw[engblue] (\cx-0.7, 2.85) -- (\cx-0.7, 3.15);
\draw[engblue] (\cx+0.7, 2.85) -- (\cx+0.7, 3.15);
\node[right, font=\scriptsize] at (\cx+0.9, 3)
  {CI for the mean $\bar{x} \pm t\,s/\sqrt{n}$};

% Prediction interval (wider), y = 2
\draw[very thick, engamber] (\cx-2.4, 2) -- (\cx+2.4, 2);
\draw[engamber] (\cx-2.4, 1.85) -- (\cx-2.4, 2.15);
\draw[engamber] (\cx+2.4, 1.85) -- (\cx+2.4, 2.15);
\node[right, font=\scriptsize] at (\cx+2.6, 2)
  {PI for next unit $\bar{x} \pm t\,s\sqrt{1+1/n}$};

% Tolerance interval (widest), y = 1
\draw[very thick, engred] (\cx-3.4, 1) -- (\cx+3.4, 1);
\draw[engred] (\cx-3.4, 0.85) -- (\cx-3.4, 1.15);
\draw[engred] (\cx+3.4, 0.85) -- (\cx+3.4, 1.15);
\node[right, font=\scriptsize] at (\cx+3.6, 1)
  {TI for $99\%$ of units $\bar{x} \pm k\,s$};

% Axis
\draw[->, thick] (0.2, 0.2) -- (9.8, 0.2)
  node[below, font=\scriptsize] {measured value};

\end{tikzpicture}
\caption[Confidence, prediction, and tolerance intervals compared]{The
three intervals of section 8.4, drawn to relative scale for a sample
of $n = 30$ from a Gaussian population at $95\,\%$ confidence. All
three share the sample mean as their centre. The confidence interval
(blue) is the narrowest: it bounds only the population mean, and its
half-width falls as $1/\sqrt{n}$. The prediction interval (amber)
adds the scatter of a single new draw and stays wide as $n$ grows.
The tolerance interval (red) is widest: it must contain a stated
fraction of the whole population with stated confidence, and it
carries both the population spread and the sampling uncertainty in
its factor $k$. A report that quotes the confidence interval where
the reader needs the prediction or tolerance interval understates
the operational scatter by the ratio of the bar lengths.}
\label{fig:vol01:ch08:interval-ladder}
\end{figure}
